\documentclass[11pt,a4paper]{article}
\usepackage[utf8]{inputenc}
\usepackage[T5]{fontenc}
\usepackage[vietnamese]{babel}
\usepackage{amsmath,amssymb,amsfonts,amsthm}
\usepackage{geometry}
\geometry{left=1.8cm,right=1.8cm,top=2cm,bottom=2cm}
\usepackage{booktabs}
\usepackage{xcolor}
\usepackage{listings}
\usepackage{fancyhdr}
\usepackage{tcolorbox}
\usepackage{array}
\usepackage{multicol}
\usepackage{hyperref}

\definecolor{primary}{rgb}{0.0, 0.25, 0.45}
\definecolor{secondary}{rgb}{0.1, 0.5, 0.3}
\definecolor{accent}{rgb}{0.75, 0.1, 0.1}
\definecolor{codebg}{rgb}{0.96, 0.96, 0.96}
\definecolor{grayborder}{rgb}{0.8, 0.8, 0.8}

\hypersetup{
    colorlinks=true,
    linkcolor=primary,
    urlcolor=primary
}

\lstset{
  language=C++,
  basicstyle=\ttfamily\small,
  keywordstyle=\color{primary}\bfseries,
  commentstyle=\color{secondary}\itshape,
  stringstyle=\color{accent},
  backgroundcolor=\color{codebg},
  frame=single,
  rulecolor=\color{grayborder},
  breaklines=true,
  tabsize=2,
  showstringspaces=false
}

\tcbuselibrary{skins,breakable}
\newtcolorbox{trapbox}[1][]{
  colback=red!5!white,
  colframe=accent,
  fonttitle=\bfseries\color{white},
  title={⚠️ BẪY ĐIỂM ÂM \& LỖI SAI THƯỜNG GẶP},
  enhanced,
  breakable
}

\newtcolorbox{corebox}[1][]{
  colback=blue!5!white,
  colframe=primary,
  fonttitle=\bfseries\color{white},
  title={💡 KIẾN THỨC CỐT LÕI \& PHƯƠNG PHÁP},
  enhanced,
  breakable
}

\pagestyle{fancy}
\fancyhf{}
\lhead{\textbf{\color{primary}CS112 \& DSA: GIÁO ÁN \& CẨM NANG CỐT LÕI}}
\rhead{\color{primary}Trường ĐH Công nghệ Thông tin - ĐHQG-HCM}
\cfoot{\thepage}

\begin{document}

\begin{center}
    {\Huge \bfseries \color{primary} Giáo án Toàn diện \& Cẩm nang Kiến thức}\[0.3cm]
    {\Large \bfseries Cấu trúc Dữ liệu, Giải thuật \& CS112}\[0.2cm]
    {\small \textbf{Giảng viên phụ trách:} Huỳnh Thị Thanh Thương (\texttt{thuonghtt@uit.edu.vn})}\[0.1cm]
    {\small \textit{Khoa Khoa học Máy tính -- Đại học Công nghệ Thông tin, ĐHQG-HCM}}\[0.2cm]
    \rule{\linewidth}{1.5pt}
\end{center}

\tableofcontents
\newpage

% ==============================================================================
% PHẦN 1
% ==============================================================================
\section{PHẦN 1: TOÁN HỌC NỀN TẢNG \& PHÉP ĐẾM SƠ CẤP CHO MÁY TÍNH (L00)}

\subsection{1.1 Bản chất toán học của máy tính và Phép đếm CPU}
Trong phân tích thuật toán, máy tính không đo thời gian tuyệt đối bằng miligiây vì phụ thuộc phần cứng, mà đo \textbf{số lượng phép toán cơ bản (Basic Operations)} khi kích thước dữ liệu $ tăng lên.

\begin{corebox}
Giả sử vòng lặp thực thi thân vòng lặp đúng $ lần cho các giá trị hợp lệ. Để thoát khỏi vòng lặp, CPU bắt buộc phải thực hiện thêm \textbf{1 phép so sánh điều kiện cuối cùng} trả về \texttt{FALSE}.
11\text{Số phép So sánh điều kiện} = \text{Số lần Lặp thực tế} + 111
\end{corebox}

\subsection{1.2 Hàm nguyên Sàn/Trần và Logarithm}
\begin{itemize}
    \item \textbf{Hàm sàn $\lfloor x \rfloor$ (Floor)}: Số nguyên lớn nhất $\le x$. Trong lập trình C++, phép chia số nguyên \texttt{int m = (l + r) / 2} chính là phép lấy hàm sàn $\lfloor \frac{l+r}{2} \rfloor$.
    \item \textbf{Hàm trần $\lceil x \rceil$ (Ceiling)}: Số nguyên nhỏ nhất $\ge x$. Dùng để xác định mốc biên tối thiểu để các vòng lặp kích hoạt.
    \item \textbf{Logarithm cơ số 2 ($\log_2 n$)}: Biểu diễn số lần ta bẻ đôi dữ liệu kích thước $ liên tiếp để đưa về $. Đẳng thức hoán đổi cơ số bắt buộc thuộc lòng: ^{\log_b c} = c^{\log_b a} \implies 2^{\log_2 n} = n$.
\end{itemize}

\subsection{1.3 Khử 5 Chuỗi Sigma ($\sum$) Phòng thi}
\begin{enumerate}
    \item \textbf{Tổng hằng số $}: $\sum_{i=L}^{U} c = c \cdot (U - L + 1)$.
    \item \textbf{Cấp số cộng}: $\sum_{i=1}^{n} i = 1 + 2 + \dots + n = \frac{n(n+1)}{2} \approx \frac{1}{2}n^2$.
    \item \textbf{Tổng bình phương}: $\sum_{i=1}^{n} i^2 = \frac{n(n+1)(2n+1)}{6} \approx \frac{1}{3}n^3$.
    \item \textbf{Cấp số nhân (công bội 2)}: $\sum_{i=0}^{n} 2^i = 2^0 + 2^1 + \dots + 2^n = 2^{n+1} - 1$.
    \item \textbf{Tổng Harmonic \& Logarithm}: $\sum_{i=1}^{n} \frac{1}{i} \approx \ln n + \gamma$; $\sum_{i=1}^{n} \lg i = \lg(n!) \approx n \lg n$.
\end{enumerate}

\begin{trapbox}
Xét vòng lặp lồng chạy ngược: \texttt{for (y=1..n) \{ for (j = n-y .. y) \{ ... \} \}}.
\begin{itemize}
    \item \textbf{Biện luận}: Vòng trong  = n-y \dots y$ chỉ chạy khi -y \le y \iff 2y \ge n \iff y \ge \lceil n/2 \rceil$.
    \item \textbf{Bẫy chết người}: Tính vội $\sum_{y=1}^n (2y - n + 2)$ sẽ sinh ra số lần lặp âm khi  < \lceil n/2 \rceil
    \item \textbf{Khắc phục chuẩn mực}: Tách tổng cô lập $+1$ phép so sánh kiểm tra ranh giới:
    11\text{So sánh} = \sum_{y=\lceil n/2 \rceil}^{n} (2y - n + 1) + \sum_{y=1}^{n} 1 = \frac{n^2 + 8n + 4}{4} \quad (\text{với } n \text{ chẵn})11
\end{itemize}
\end{trapbox}

\newpage
% ==============================================================================
% PHẦN 2
% ==============================================================================
\section{PHẦN 2: CẤU TRÚC DỮ LIỆU CỐT LÕI -- NỀN TẢNG DSA (L01--L05)}

\subsection{2.1 Mảng tĩnh, Mảng động \& Danh sách liên kết (L01)}
\begin{table}[h]
\centering
\caption{So sánh đặc tính hiệu năng giữa Mảng và Danh sách liên kết}
\vspace{0.2cm}
\begin{tabular}{p{4.5cm} p{5.5cm} p{5.5cm}}
\toprule
\textbf{Thao tác} & \textbf{Mảng (Array / Vector)} & \textbf{Danh sách liên kết (Linked List)} \\
\midrule
Truy xuất theo chỉ số $ & (1)$ (Địa chỉ $\text{Base} + i \cdot \text{size}$) & (n)$ (Duyệt con trỏ từ Head) \\
Chèn / Xóa ở đầu & (n)$ (Phải dịch chuyển các phần tử) & (1)$ (Thay đổi con trỏ Head) \\
Chèn / Xóa ở cuối & (1)$ amortized (Mảng động) & (1)$ (Nếu có con trỏ Tail) \\
Chèn / Xóa ở giữa & (n)$ & (1)$ (Khi đã có con trỏ nút kề) \\
Bộ nhớ lưu trữ & Ô nhớ liên tục, tối ưu cache CPU & Ô nhớ phân tán, tốn thêm con trỏ \\
\bottomrule
\end{tabular}
\end{table}

\subsection{2.2 Ngăn xếp (Stack) \& Hàng đợi (Queue) (L02)}
\begin{itemize}
    \item \textbf{Stack (LIFO - Last In First Out)}: Thao tác \texttt{Push}, \texttt{Pop}, \texttt{Top} đều tốn (1)$. Ứng dụng quản lý Call Stack hệ điều hành khi gọi đệ quy, kiểm tra chuỗi ngoặc hợp lệ, chuyển đổi biểu thức Infix sang Postfix.
    \item \textbf{Queue (FIFO - First In First Out)}: Thao tác \texttt{Enqueue}, \texttt{Dequeue} đều tốn (1)$. Hàng đợi vòng (Circular Queue) sử dụng phép toán chia lấy dư: \texttt{rear = (rear + 1) \% MAX}.
\end{itemize}

\subsection{2.3 Cây nhị phân, BST \& Cây AVL cân bằng (L03)}
\begin{itemize}
    \item \textbf{Cây tìm kiếm nhị phân (BST)}: Tính chất Trái $<$ Gốc $<$ Phải. Phép duyệt \textbf{In-order (LNR)} trên BST luôn thu được dãy phần tử tăng dần có thứ tự.
    \item \textbf{Cây AVL}: Đảm bảo chênh lệch chiều cao giữa 2 cây con (Hệ số cân bằng $\text{BF} = h_{\text{Left}} - h_{\text{Right}}$) luôn nằm trong $\{-1, 0, 1\}$. Sử dụng 4 phép quay cân bằng: Quay đơn (LL, RR) và Quay đôi (LR, RL) để duy trì độ phức tạp tìm kiếm (\log n)$.
\end{itemize}

\subsection{2.4 Bảng băm (Hash Table) \& Heap / Priority Queue (L04)}
\begin{itemize}
    \item \textbf{Bảng băm}: Ánh xạ khóa $ thành chỉ số mảng qua hàm băm (k) = k \bmod m$. Giải quyết đụng độ bằng \textbf{Chaining} (Danh sách liên kết tại mỗi bucket) hoặc \textbf{Open Addressing} (Linear Probing). Chi phí tra cứu (1)$ trung bình, worst-case (n)$.
    \item \textbf{Binary Heap}: Cây nhị phân gần hoàn chỉnh lưu trên mảng 1D. Nút cha tại $, con trái tại i+1$, con phải tại i+2$. Thao tác \texttt{Build-Heap} tốn (n)$, \texttt{Heapify} tốn (\log n)$.
\end{itemize}

\subsection{2.5 Cấu trúc Đồ thị \& Thuật toán Duyệt BFS, DFS (L05)}
\begin{itemize}
    \item \textbf{Ma trận kề} ($|V| \times |V|$): Tốn bộ nhớ (V^2)$, kiểm tra kề (1)$. Phù hợp đồ thị dày.
    \item \textbf{Danh sách kề}: Tốn bộ nhớ (V + E)$. Phù hợp đồ thị thưa.
    \item \textbf{Duyệt BFS (dùng Queue)}: Tìm đường đi ngắn nhất trên đồ thị không trọng số, độ phức tạp (V + E)$.
    \item \textbf{Duyệt DFS (dùng Stack / Đệ quy)}: Kiểm tra tính liên thông, tìm chu trình, sắp xếp Topo, độ phức tạp (V + E)$.
\end{itemize}

\newpage
% ==============================================================================
% PHẦN 3
% ==============================================================================
\section{PHẦN 3: PHÂN TÍCH HIỆU NĂNG TIỆM CẬN (L06--L07)}

\subsection{3.1 5 Ký hiệu Tiệm cận Hình thức \& Phương pháp Giới hạn (L06)}
\begin{enumerate}
    \item \textbf{Chặn trên Big-O ($)}: (n) \in O(g(n)) \iff \exists c > 0, n_0 > 0 : 0 \le f(n) \le c \cdot g(n), \quad \forall n \ge n_0$.
    \item \textbf{Chặn dưới Big-Omega ($\Omega$)}: (n) \in \Omega(g(n)) \iff \exists c > 0, n_0 > 0 : 0 \le c \cdot g(n) \le f(n), \quad \forall n \ge n_0$.
    \item \textbf{Chặn chặt Big-Theta ($\Theta$)}: (n) \in \Theta(g(n)) \iff \exists c_1 > 0, c_2 > 0, n_0 > 0 : c_1 g(n) \le f(n) \le c_2 g(n), \quad \forall n \ge n_0$.
    \item \textbf{Little-o ($)}: $\lim_{n \to \infty} \frac{f(n)}{g(n)} = 0 \implies f(n) \in o(g(n))$.
    \item \textbf{Little-omega ($\omega$)}: $\lim_{n \to \infty} \frac{f(n)}{g(n)} = \infty \implies f(n) \in \omega(g(n))$.
\end{enumerate}

\subsection{3.2 Chuỗi 15 Bậc Thứ tự Tăng trưởng Tiệm cận (L07)}
Bắt buộc thuộc lòng chuỗi thứ tự tăng dần chuẩn mực:
111 \prec \log\log n \prec \log n \prec \log^2 n \prec \sqrt[3]{n} \prec \sqrt{n} \prec n \prec n\log n \prec n^2 \prec n^3 \prec n^4 \prec 2^n \prec 3^n \prec 4^n \prec n! \prec n^n11

\begin{trapbox}
\begin{itemize}
    \item \textbf{Sai lầm}: Nhầm ^{2n}$ đồng bậc với ^n$.
    \item \textbf{Sửa chuẩn đại số}: ^{2n} = (2^2)^n = 4^n$.
    \item \textbf{Chứng minh giới hạn}: $\lim_{n \to \infty} \frac{2^{2n}}{2^n} = \lim_{n \to \infty} \frac{4^n}{2^n} = \lim_{n \to \infty} 2^n = \infty \implies 2^{2n} \in \omega(2^n)$. Do đó ^{2n} \succ 3^n \succ 2^n$.
\end{itemize}
\end{trapbox}

\newpage
% ==============================================================================
% PHẦN 4
% ==============================================================================
\section{PHẦN 4: PHÂN TÍCH GIẢI THUẬT ĐỆ QUY (L08--L11)}

\subsection{4.1 Phương pháp Thế ngược -- Bài toán Tháp Hà Nội (L08)}
Hệ thức: (1) = 1$; (n) = 2T(n-1) + 1 \quad (\forall n > 1)$.
\begin{enumerate}
    \item Thế lượt 1: (n) = 2[2T(n-2) + 1] + 1 = 2^2 T(n-2) + 2^1 + 2^0$.
    \item Thế lượt $: (n) = 2^i T(n-i) + \sum_{j=0}^{i-1} 2^j = 2^i T(n-i) + (2^i - 1)$.
    \item Mốc dừng  - i = 1 \iff i = n - 1$.
    \item Thế  = n - 1$: (n) = 2^{n-1} T(1) + (2^{n-1} - 1) = \mathbf{2^n - 1 \in \Theta(2^n)}$.
\end{enumerate}

\subsection{4.2 Phương pháp Hàm sinh (Generating Function) (L09)}
Định nghĩa (x) = \sum_{n=1}^{\infty} T(n) x^n$.
Nhân ^n$ vào hệ thức (n) = 2T(n-1) + 1$ và lấy tổng từ  = 2 \to \infty$:
11G(x) - x = 2x G(x) + \frac{x^2}{1-x} \implies G(x)(1-2x) = \frac{x}{1-x} \implies G(x) = \frac{x}{(1-x)(1-2x)}11
Phân rã phân thức tối giản:
11G(x) = \frac{1}{1-2x} - \frac{1}{1-x} = \sum_{n=1}^{\infty} 2^n x^n - \sum_{n=1}^{\infty} x^n = \sum_{n=1}^{\infty} (2^n - 1) x^n \implies \mathbf{T(n) = 2^n - 1}11

\subsection{4.3 Thuật toán findMax \& Bẫy số nhánh $ (L10)}
\begin{lstlisting}
int findMax(int a[], int l, int r) {
    if (l == r) return a[l];
    int m = (l + r) / 2;
    if (a[m] < a[m+1]) return findMax(a, m+1, r);
    else if (a[m] > a[m+1]) return findMax(a, l, m);
    else return max(findMax(a, m+1, r), findMax(a, l, m));
}
\end{lstlisting}
\begin{itemize}
    \item \textbf{Best/Average case}: Do nằm trong cấu trúc \texttt{if..else if..} loại trừ, CPU chỉ rẽ vào 1 nhánh đệ quy $\implies \mathbf{a = 1} \implies T(n) = T(n/2) + 1 \implies \mathbf{\Theta(\log n)}$.
    \item \textbf{Worst case} (các phần tử bằng nhau rơi vào \texttt{else}): CPU chạy cả 2 nhánh $\implies \mathbf{a = 2} \implies T(n) = 2T(n/2) + 2 \implies \mathbf{\Theta(n)}$.
    \item \textbf{Chứng minh quy nạp (n) = 2n - 2$}: Cơ sở =1 \implies T(1)=0$ đúng. Giả thiết (k) = 2k-2 \quad (\forall k < n)$. Bước quy nạp: (n) = 2 T(n/2) + 2 = 2(2(n/2) - 2) + 2 = 2n - 2$ (Q.E.D).
\end{itemize}

\begin{trapbox}
Khi chứng minh (n) \le f(n)$, sau khi biến đổi ra (n) \le f(n) + R(n)$, điều kiện tiên quyết bắt buộc phải có là $\mathbf{R(n) \le 0}$. Nếu (n) > 0$, chiều bất đẳng thức bị vỡ và phép chứng minh sai hoàn toàn!
\end{trapbox}

\subsection{4.4 Kỹ thuật Đổi biến số giải hệ thức căn thức (L11)}
Giải hệ thức (n) = \sqrt{n} T(\sqrt{n}) + n$:
\begin{enumerate}
    \item Đặt  = 2^m \iff m = \log_2 n \implies \sqrt{n} = 2^{m/2}$.
    \item Thay vào phương trình: (2^m) = 2^{m/2} T(2^{m/2}) + 2^m$.
    \item Đặt (m) = T(2^m) \implies S(m) = 2^{m/2} S(m/2) + 2^m$.
    \item Chia hai vế cho ^m$: $\frac{S(m)}{2^m} = \frac{S(m/2)}{2^{m/2}} + 1$.
    \item Đặt (m) = \frac{S(m)}{2^m} \implies R(m) = R(m/2) + 1 \implies R(m) = \log_2 m + c$.
    \item Thế ngược (m) = 2^m (\log_2 m + c)$.
    \item Thế ngược biến gốc: (n) = n \log_2(\log_2 n) + c \cdot n \implies \mathbf{T(n) \in \Theta(n \log \log n)}$.
\end{enumerate}

\newpage
% ==============================================================================
% PHẦN 5
% ==============================================================================
\section{PHẦN 5: CÁC CHIẾN LƯỢC THIẾT KẾ THUẬT TOÁN (L12--L17)}

\subsection{5.1 Chiến lược Chia để trị \& Định lý Master (L12)}
Phương trình tổng quát: (n) = a T(n/b) + f(n)$. So sánh (n)$ với ^{\log_b a}$:
\begin{itemize}
    \item \textbf{Case 1}: (n) \in O(n^{\log_b a - \epsilon}) \implies T(n) \in \Theta(n^{\log_b a})$.
    \item \textbf{Case 2}: (n) \in \Theta(n^{\log_b a}) \implies T(n) \in \Theta(n^{\log_b a} \log n)$.
    \item \textbf{Case 3}: (n) \in \Omega(n^{\log_b a + \epsilon})$ và  f(n/b) \le c f(n) \implies T(n) \in \Theta(f(n))$.
    \item \textbf{Ví dụ Mergesort}: (n) = 2T(n/2) + n \implies a=2, b=2 \implies n^{\log_2 2} = n^1 \implies \mathbf{T(n) \in \Theta(n \log n)}$.
\end{itemize}

\subsection{5.2 Quay lui (Backtracking) \& Phục hồi trạng thái (L14)}
\begin{lstlisting}
void Backtrack(int i) {
    for (int j = 1; j <= m; j++) {
        if (isValid(i, j)) {
            x[i] = a[j];
            visited[j] = true;      // Mark state
            if (i == n) printSolution();
            else Backtrack(i + 1);
            visited[j] = false;     // UNMARK STATE (PHỤC HỒI TRẠNG THÁI SỐNG CÒN)
        }
    }
}
\end{lstlisting}

\subsection{5.3 Nhánh cận -- Ma trận rút gọn TSP (L15)}
Cho ma trận chi phí $  \times 5$.
\begin{enumerate}
    \item \textbf{Rút gọn hàng}: Tìm  = \min_j C[i][j]$ cho từng hàng. Trừ hàng $ cho $. Tổng  = \sum r_i = 10+2+2+3+4 = 21$.
    \item \textbf{Rút gọn cột}: Tìm  = \min_i C'[i][j]$ cho từng cột TRÊN MA TRẬN MỚI '$. Trừ cột $ cho $. Tổng {\text{col}} = \sum c_j = 1+0+3+0+0 = 4$.
    \item \textbf{Chặn dưới nút gốc}: $\mathbf{LB_{\text{gốc}} = R + C_{\text{col}} = 21 + 4 = 25}$.
\end{enumerate}

\subsection{5.4 Tham lam -- Ba lô phân rã Fractional Knapsack (L16)}
Sức chứa  = 20$,  = (18, 15, 10)$,  = (25, 24, 15)$.
\begin{enumerate}
    \item Tính đơn giá  = p_i / w_i$:  = 1.389$,  = 1.600$,  = 1.500$.
    \item Sắp xếp ưu tiên giảm dần: Đồ vật  \to$ Đồ vật  \to$ Đồ vật $.
    \item Lấy  = 1 \implies Rem = 20 - 15 = 5$ kg,  = 24$.
    \item Lấy  = Rem / w_3 = 5 / 10 = \mathbf{0.5} \implies Rem = 0$ kg,  = 24 + 0.5(15) = \mathbf{31.5}$.
    \item Nghiệm tối ưu:  = (0, 1, 0.5)$,  = \mathbf{31.5}$.
\end{enumerate}

\subsection{5.5 Tham lam -- Lập lịch hội trường Activity Scheduler (L17)}
Sắp xếp tất cả sự kiện theo \textbf{thời gian kết thúc $ tăng dần}. Luôn chọn sự kiện kết thúc sớm nhất chừa khoảng trống cho các sự kiện sau $\implies$ Độ phức tạp $\mathbf{\Theta(n \log n)}$.

\newpage
% ==============================================================================
% PHẦN 6
% ==============================================================================
\section{PHẦN 6: QUY HOẠCH ĐỘNG (L18--L20)}

\subsection{6.1 Bản chất Quy hoạch động \& Bài toán Đặt ngoặc MCM (L19)}
Cho biểu thức  = 6 + 0 \times 6 + 0.1$ với  = (6, 0, 6, 0.1)$ và toán tử  = (+, \times, +)$.
\begin{itemize}
    \item \textbf{Cơ sở}: [i][i] = x_i \implies DP[0][0]=6, DP[1][1]=0, DP[2][2]=6, DP[3][3]=0.1$.
    \item \textbf{Công thức truy hồi}: [i][j] = \max_{i \le k < j} \{ DP[i][k] \ \mathbf{o_k} \ DP[k+1][j] \}$.
    \item \textbf{Điền theo độ dài đoạn  = j - i + 1$}:
    \begin{itemize}
        \item =2$: [0][1] = 6+0 = 6$; [1][2] = 0 \times 6 = 0$; [2][3] = 6 + 0.1 = 6.1$.
        \item =3$: [0][2] = \max(6+0, 6 \times 6) = \mathbf{36}$ (tại =1$). [1][3] = \max(0 \times 6.1, 0+0.1) = \mathbf{0.1}$ (tại =2$).
        \item =4$ (Ô đích [0][3]$): Thử =0 \implies 6+0.1=6.1$; =1 \implies 6 \times 6.1 = \mathbf{36.6}$; =2 \implies 36+0.1=36.1 \implies \mathbf{DP[0][3] = 36.6}$ tại =1$.
    \end{itemize}
\end{itemize}

\begin{table}[h]
\centering
\caption{Bảng Ma trận DP[i][j] trình bày bài thi tự luận}
\vspace{0.2cm}
\begin{tabular}{ccccc}
\toprule
  \setminus j$ & \textbf{0} & \textbf{1} & \textbf{2} & \textbf{3} \\
\midrule
\textbf{0} & \textbf{6} & \textbf{6} & \textbf{36} & \textbf{36.6} (tại =1$) \\
\textbf{1} & - & \textbf{0} & \textbf{0} & \textbf{0.1} \\
\textbf{2} & - & - & \textbf{6} & \textbf{6.1} \\
\textbf{3} & - & - & - & \textbf{0.1} \\
\bottomrule
\end{tabular}
\end{table}

\textbf{Truy vết biểu thức đóng ngoặc}: Cắt tại =1$ (toán tử $\times$) $\implies \mathbf{(6 + 0) \times (6 + 0.1) = 36.6}$.

\newpage
% ==============================================================================
% PHẦN 7
% ==============================================================================
\section{PHẦN 7: TỐI ƯU HÓA TUYẾN TÍNH \& HÌNH HỌC (L21--L22)}

\subsection{7.1 Quy hoạch Tuyến tính LP 2D -- Phương pháp Hình học (L21)}
Maximize  = 15x_1 + 10x_2$ thỏa  + x_2 \le 24$; x_1 + x_2 \le 32$;  \ge 0, x_2 \ge 0$.
\begin{enumerate}
    \item Tìm giao điểm 4 đỉnh đa giác miền chấp nhận được $: (0, 0)$, (16, 0)$, (0, 24)$,  = d_1 \cap d_2 \implies B(8, 16)$.
    \item Thế tọa độ 4 đỉnh vào $: (O)=0$, (A)=240$, (C)=240$, (B) = 15(8) + 10(16) = \mathbf{280}$.
    \item \textbf{Kết luận}: Lợi nhuận cực đại {\max} = \mathbf{280}$ USD đạt được tại đỉnh (8, 16)$.
\end{enumerate}

\newpage
% ==============================================================================
% PHẦN 8
% ==============================================================================
\section{PHẦN 8: BÀI TẬP LỚN BTL CS112 (10/10) \& CẨM NANG PHÒNG THI (L23--L24)}

\subsection{8.1 Chuẩn mực Bài Tập Lớn CS112 (L23)}
\begin{itemize}
    \item \textbf{Cấu trúc đóng gói zip}: \texttt{1\_SourceCode/}, \texttt{2\_VideoDemo/}, \texttt{3\_Testcase/} (5 nhóm), \texttt{4\_Report/} (.docx/.tex và .pdf).
    \item \textbf{Contract Bottleneck 1 (Mục 8 Báo cáo)}: Bắt buộc lập bảng \textbf{Nhật ký Prompt AI 5 cột} (Công đoạn, AI Model, Prompt gốc, Kết quả AI, Cách SV kiểm chứng và sửa chữa).
    \item \textbf{Contract Bottleneck 2 (Demo Visualizer)}: Bắt buộc cài đặt so sánh đồng thời 4 thuật toán (\textit{Bottom-up DP, Top-down Memoization, Brute-force Recursion, Greedy Split}) kèm biểu đồ Benchmark thời gian và bộ nhớ.
\end{itemize}

\subsection{8.2 Chiến thuật Phân bổ Thời gian Phòng thi 90 Phút (L24)}
\begin{table}[h]
\centering
\caption{Kế hoạch tác chiến 90 phút thi tự luận CS112 đạt điểm 10/10}
\vspace{0.2cm}
\begin{tabular}{p{3.5cm} p{9.5cm} p{2.5cm}}
\toprule
\textbf{Thời gian} & \textbf{Nội dung công việc} & \textbf{Thang điểm} \\
\midrule
00 -- 20 phút & \textbf{Câu 1}: Tiệm cận (Chứng minh $\Theta$ + Sắp xếp chuỗi 15 bậc) & 2.5 điểm \\
20 -- 45 phút & \textbf{Câu 2}: Đệ quy (Tháp Hà Nội / findMax / Đổi biến căn thức) & 2.5 điểm \\
45 -- 70 phút & \textbf{Câu 3}: Quy hoạch động Đặt ngoặc (Lập bảng $ + Truy vết) & 2.5 điểm \\
70 -- 90 phút & \textbf{Câu Tự chọn}: Tham lam Knapsack / LP 2D / Reduced Matrix TSP & 2.5 điểm \\
\bottomrule
\end{tabular}
\end{table}

\end{document}
